% EGA IV-3 standalone-English terminal matter, printed pp. 249--255.
% Source-grounded independent English edition.

\begin{center}
\emph{(To be continued.)}
\end{center}

\clearpage
\oldpage[IV-3]{249}
\phantomsection
\addcontentsline{toc}{section}{Bibliography (continued)}
\section*{BIBLIOGRAPHY (\emph{continued})}

\begingroup
\small
\setlength{\parindent}{0pt}
\setlength{\parskip}{3pt}
[39] J.-P.~\textsc{Serre}, \emph{Local Fields}, Paris (Hermann), 1962.

[40] J.-P.~\textsc{Serre}, \emph{Proalgebraic Groups}, \emph{Publ. Inst. Hautes Études Sci.}, no.~7 (1960).

\endgroup

\clearpage
\oldpage[IV-3]{250}
\phantomsection
\addcontentsline{toc}{section}{Index of notation}
\section*{INDEX OF NOTATION}

\begingroup
\small
\setlength{\parindent}{0pt}
\setlength{\parskip}{2pt}

$\mathfrak C(X)$, $\mathfrak D_c(X)$, $\mathfrak F_c(X)$,
$\mathfrak{LF}_c(X)$ ($X$ a prescheme) : 8.3.9.

$\mathfrak Q(\mathcal F)$ ($\mathcal F$ an $\mathcal O_X$-module) : 8.5.10.

$\mathfrak{Spr}(Y)$, $\mathfrak{Spro}(Y)$, $\mathfrak{Sprf}(Y)$
($Y$ a prescheme) : 8.6.1.

$\mathbf{Pro}(\mathcal C)$ ($\mathcal C$ a category) : 8.13.3.

$X_s$, $\mathcal F_s$, $u_s$, $g_s$, $Z_s$, $h_s$
($X$ an $S$-prescheme, $s\in S$, $\mathcal F$ an $\mathcal O_X$-module,
$u$ a homomorphism of $\mathcal O_X$-modules, $g$ a section of an
$\mathcal O_X$-module, $Z$ a subset of $X$, $h:X\to Y$ an $S$-morphism) : 9.4.1.

$\mathfrak{Qc}(X)$, $\mathfrak{Lqc}(X)$, $\mathfrak O(X)$,
$\mathfrak F(X)$, $\mathfrak{Lf}(X)$, $\mathfrak{Lc}(X)$
($X$ a topological space) : 10.1.1.

$\operatorname{prof}^{*}_{x}(\mathcal F)$,
$\operatorname{prof}^{*}_{Z}(\mathcal F)$,
$\operatorname{prof}^{*}(\mathcal F)$
($\mathcal F$ an $\mathcal O_X$-module) : 10.8.1.

$S(X)$ ($X$ a Jacobson prescheme) : 10.9.1.

$S(f)$ ($f$ a morphism of Jacobson preschemes) : 10.9.2.

$\operatorname{Spm}(A)$, $D^m(h)$ ($A$ a Jacobson ring, $h\in A$) : 10.9.3.

$\operatorname{Tor.dim}_B(M)$ ($M$ a $B$-module) : 12.3.2.

\endgroup

\clearpage
\oldpage[IV-3]{251}
\phantomsection
\addcontentsline{toc}{section}{Terminological index}
\section*{TERMINOLOGICAL INDEX}

\begingroup
\footnotesize
\setlength{\parindent}{0pt}
\setlength{\parskip}{1.1pt}
\newcommand{\EGAIVthreeTerm}[2]{#1 : #2.\par}

\EGAIVthreeTerm{Flatness criterion by fibres}{11.3.10}
\EGAIVthreeTerm{Algebraic space over $k$, $k$-algebraic space, prealgebraic space over $k$, $k$-prealgebraic space}{10.10.2}
\EGAIVthreeTerm{Jacobson space}{10.3.1}
\EGAIVthreeTerm{Reduced $k$-prealgebraic space}{10.10.4}
\EGAIVthreeTerm{Schematically dominant family of morphisms}{11.10.2}
\EGAIVthreeTerm{Universally schematically dominant family of morphisms}{11.10.7}
\EGAIVthreeTerm{Schematically dense family of subpreschemes}{11.10.2}
\EGAIVthreeTerm{Separating family of homomorphisms of sheaves of modules}{11.9.1}
\EGAIVthreeTerm{Separating family of module homomorphisms}{11.9.4}
\EGAIVthreeTerm{Universally separating family of homomorphisms of sheaves of modules}{11.9.14}
\EGAIVthreeTerm{Geometrically isomorphic ($k$-preschemes)}{9.1.4}
\EGAIVthreeTerm{Universally injective homomorphism of modules}{11.9.14 and 11.9.18}
\EGAIVthreeTerm{Chow's lemma for morphisms of finite presentation}{8.10.5.1}
\EGAIVthreeTerm{Zariski's ``Main Theorem'' for morphisms of finite presentation}{8.12.6}
\EGAIVthreeTerm{Equidimensional morphism at a point, equidimensional morphism}{13.2.2 and 13.3.2}
\EGAIVthreeTerm{Morphism proper at a point}{15.7.1}
\EGAIVthreeTerm{Pseudo-finite morphism}{8.12.3}
\EGAIVthreeTerm{Submersive morphism, universally submersive morphism}{15.7.8}
\EGAIVthreeTerm{Universally open morphism at a point}{14.3.3}
\EGAIVthreeTerm{Ultra-affine open subset}{10.9.5}
\EGAIVthreeTerm{Saturated finite subset of an algebraic $k$-prescheme}{9.8.9}
\EGAIVthreeTerm{Subset proper at a point (of an $S$-prescheme)}{15.7.1}
\EGAIVthreeTerm{Quasi-constructible subset, locally quasi-constructible subset of a topological space}{10.1.1}
\EGAIVthreeTerm{Very dense subset of a topological space}{10.1.3}
\EGAIVthreeTerm{Geometrically irreducible polynomial}{9.7.4}
\EGAIVthreeTerm{Jacobson prescheme}{10.4.1}
\EGAIVthreeTerm{Prescheme equidimensional over another at a point, equidimensional over another}{13.2.2 and 13.3.2}
\EGAIVthreeTerm{Prescheme essentially affine over another}{8.13.4}
\EGAIVthreeTerm{Finite extension principle}{9.1.1}
\EGAIVthreeTerm{Rectified depth}{10.8.1}
\EGAIVthreeTerm{Pro-object, morphism of pro-objects}{8.13.3}
\EGAIVthreeTerm{Essentially affine pro-object}{8.13.4, 8.14.1}
\EGAIVthreeTerm{Constructible property, ind-constructible property}{9.2.1 and 9.2.2, (i) and (ii)}
\EGAIVthreeTerm{Quasi-homeomorphism of topological spaces}{10.2.2}
\EGAIVthreeTerm{Quasi-isomorphism of ringed spaces}{10.2.8}
\EGAIVthreeTerm{Saturate of a finite subset of an algebraic $k$-prescheme}{9.8.9}
\EGAIVthreeTerm{Subset of an $S$-prescheme submersive, universally submersive over its image}{15.7.8}
\EGAIVthreeTerm{Primary skeleton of a module, virtual skeleton}{9.8.9}
\EGAIVthreeTerm{Maximal spectrum of a Jacobson ring}{10.9.3}
\EGAIVthreeTerm{Primary type of a module, geometric primary type of a module}{9.8.9}
\EGAIVthreeTerm{Ultra-prescheme, morphism of ultra-preschemes}{10.9.5}
\EGAIVthreeTerm{$R$-ultra-prescheme}{10.10.2}

\let\EGAIVthreeTerm\relax
\endgroup

\clearpage
% Printed page 252 is blank.
\thispagestyle{empty}
\null

\clearpage
\oldpage[IV-3]{253}
\phantomsection
\addcontentsline{toc}{section}{Table of contents}
\section*{TABLE OF CONTENTS}

\begingroup
\small
\setlength{\parindent}{0pt}
\setlength{\parskip}{1.5pt}
\newcommand{\EGAIVthreeContentsLine}[3]{\textbf{#1}\quad #2\unskip\nobreak\hspace{0.4em}\nobreak\dotfill\nobreak\mbox{#3}\par}

\EGAIVthreeContentsLine{CHAPTER IV.}{\textbf{Local study of schemes and morphisms of schemes} \emph{(continued)}}{5}
\EGAIVthreeContentsLine{\S~8.}{Projective limits of preschemes}{5}
\EGAIVthreeContentsLine{8.1.}{Introduction}{5}
\EGAIVthreeContentsLine{8.2.}{Projective limits of preschemes}{7}
\EGAIVthreeContentsLine{8.3.}{Constructible subsets in a projective limit of preschemes}{12}
\EGAIVthreeContentsLine{8.4.}{Irreducibility and connectedness criteria for projective limits of preschemes}{17}
\EGAIVthreeContentsLine{8.5.}{Modules of finite presentation on a projective limit of preschemes}{19}
\EGAIVthreeContentsLine{8.6.}{Subpreschemes of finite presentation of a projective limit of preschemes}{25}
\EGAIVthreeContentsLine{8.7.}{Criteria for a projective limit of preschemes to be a reduced (resp. integral) prescheme}{27}
\EGAIVthreeContentsLine{8.8.}{Preschemes of finite presentation on a projective limit of preschemes}{28}
\EGAIVthreeContentsLine{8.9.}{First applications to the elimination of Noetherian hypotheses}{34}
\EGAIVthreeContentsLine{8.10.}{Stability properties of morphisms under passage to the projective limit}{36}
\EGAIVthreeContentsLine{8.11.}{Application to quasi-finite morphisms}{41}
\EGAIVthreeContentsLine{8.12.}{New proof and generalization of Zariski's ``Main Theorem''}{43}
\EGAIVthreeContentsLine{8.13.}{Translation in terms of pro-objects}{49}
\EGAIVthreeContentsLine{8.14.}{Characterization of a prescheme locally of finite presentation over another, in terms of the functor it represents}{52}

\EGAIVthreeContentsLine{\S~9.}{Constructible properties}{54}
\EGAIVthreeContentsLine{9.1.}{The finite extension principle}{54}
\EGAIVthreeContentsLine{9.2.}{Constructible and ind-constructible properties}{56}
\EGAIVthreeContentsLine{9.3.}{Constructible properties of morphisms of algebraic preschemes}{60}

\clearpage
\oldpage[IV-3]{254}
\EGAIVthreeContentsLine{9.4.}{Constructibility of certain properties of modules}{62}
\EGAIVthreeContentsLine{9.5.}{Constructibility of topological properties}{67}
\EGAIVthreeContentsLine{9.6.}{Constructibility of certain properties of morphisms}{71}
\EGAIVthreeContentsLine{9.7.}{Constructibility of separability, geometric irreducibility, and geometric connectedness}{76}
\EGAIVthreeContentsLine{9.8.}{Primary decomposition near a generic fibre}{83}
\EGAIVthreeContentsLine{9.9.}{Constructibility of local properties of fibres}{88}

\EGAIVthreeContentsLine{\S~10.}{Jacobson preschemes}{95}
\EGAIVthreeContentsLine{10.1.}{Very dense subsets of a topological space}{95}
\EGAIVthreeContentsLine{10.2.}{Quasi-homeomorphisms}{97}
\EGAIVthreeContentsLine{10.3.}{Jacobson spaces}{101}
\EGAIVthreeContentsLine{10.4.}{Jacobson preschemes and Jacobson rings}{101}
\EGAIVthreeContentsLine{10.5.}{Noetherian Jacobson preschemes}{104}
\EGAIVthreeContentsLine{10.6.}{Dimension in Jacobson preschemes}{107}
\EGAIVthreeContentsLine{10.7.}{Examples and counterexamples}{109}
\EGAIVthreeContentsLine{10.8.}{Rectified depth}{110}
\EGAIVthreeContentsLine{10.9.}{Maximal spectra and ultrapreschemes}{112}
\EGAIVthreeContentsLine{10.10.}{Serre algebraic spaces}{114}

\EGAIVthreeContentsLine{\S~11.}{Topological properties of flat morphisms of finite presentation. Flatness criteria}{116}
\EGAIVthreeContentsLine{11.1.}{Flatness loci (Noetherian case)}{117}
\EGAIVthreeContentsLine{11.2.}{Flatness of a projective limit of preschemes}{119}
\EGAIVthreeContentsLine{11.3.}{Application to the elimination of Noetherian hypotheses}{132}
\EGAIVthreeContentsLine{11.4.}{Descent of flatness by arbitrary morphisms: case of an Artinian base prescheme}{143}
\EGAIVthreeContentsLine{11.5.}{Descent of flatness by arbitrary morphisms: general case}{150}
\EGAIVthreeContentsLine{11.6.}{Descent of flatness by arbitrary morphisms: case of a unibranch base prescheme}{154}
\EGAIVthreeContentsLine{11.7.}{Counterexamples}{157}
\EGAIVthreeContentsLine{11.8.}{A valuative criterion for flatness}{159}
\EGAIVthreeContentsLine{11.9.}{Separating and universally separating families of homomorphisms of sheaves of modules}{160}
\EGAIVthreeContentsLine{11.10.}{Schematically dominant families of morphisms and schematically dense families of subpreschemes}{170}

\EGAIVthreeContentsLine{\S~12.}{Study of the fibres of flat morphisms of finite presentation}{173}
\EGAIVthreeContentsLine{12.0.}{Introduction}{173}
\EGAIVthreeContentsLine{12.1.}{Local properties of fibres of a flat morphism locally of finite presentation}{174}

\clearpage
\oldpage[IV-3]{255}
\EGAIVthreeContentsLine{12.2.}{Local and global properties of fibres of a proper, flat morphism of finite presentation}{179}
\EGAIVthreeContentsLine{12.3.}{Local cohomological properties of fibres of a flat morphism locally of finite presentation}{183}

\EGAIVthreeContentsLine{\S~13.}{Equidimensional morphisms}{187}
\EGAIVthreeContentsLine{13.1.}{Chevalley's semicontinuity theorem}{188}
\EGAIVthreeContentsLine{13.2.}{Equidimensional morphisms: case of dominant morphisms of irreducible preschemes}{190}
\EGAIVthreeContentsLine{13.3.}{Equidimensional morphisms: general case}{194}

\EGAIVthreeContentsLine{\S~14.}{Universally open morphisms}{199}
\EGAIVthreeContentsLine{14.1.}{Open morphisms}{200}
\EGAIVthreeContentsLine{14.2.}{Open morphisms and the dimension formula}{202}
\EGAIVthreeContentsLine{14.3.}{Universally open morphisms}{204}
\EGAIVthreeContentsLine{14.4.}{Chevalley's criterion for universally open morphisms}{209}
\EGAIVthreeContentsLine{14.5.}{Universally open morphisms and quasi-sections}{216}

\EGAIVthreeContentsLine{\S~15.}{Study of the fibres of a universally open morphism}{223}
\EGAIVthreeContentsLine{15.1.}{Multiplicities of the fibres of a universally open morphism}{223}
\EGAIVthreeContentsLine{15.2.}{Flatness of universally open morphisms with geometrically reduced fibres}{226}
\EGAIVthreeContentsLine{15.3.}{Application: criteria of reducedness and irreducibility}{228}
\EGAIVthreeContentsLine{15.4.}{Complements on Cohen-Macaulay morphisms}{229}
\EGAIVthreeContentsLine{15.5.}{Separable rank of the fibres of a quasi-finite and universally open morphism. Application to the geometric connected components of the fibres of a proper morphism}{231}
\EGAIVthreeContentsLine{15.6.}{Connected components of fibres along a section}{236}
\EGAIVthreeContentsLine{15.7.}{Appendix: Valuative criteria of local properness}{242}
\EGAIVthreeContentsLine{}{BIBLIOGRAPHY \emph{(continued)}}{249}
\EGAIVthreeContentsLine{}{INDEX OF NOTATION}{250}
\EGAIVthreeContentsLine{}{TERMINOLOGICAL INDEX}{251}

\vspace{2em}
\begin{center}
\emph{Received November 15, 1964.}
\end{center}

\let\EGAIVthreeContentsLine\relax
\endgroup
